% ╔══════════════════════════════════════════════════════════╗
% ║  CV – İlan Eşleşmesi: Final Ödev-2 Raporu               ║
% ║  Doğal Dil İşleme | Gümüşhane Üniversitesi              ║
% ║  Derleyici: pdflatex (Overleaf)                         ║
% ╚══════════════════════════════════════════════════════════╝

\documentclass[12pt, a4paper]{article}

\usepackage[utf8]{inputenc}
\usepackage[T1]{fontenc}
\usepackage[turkish]{babel}
\usepackage{geometry}
\usepackage{amsmath}
\usepackage{amssymb}
\usepackage{graphicx}
\usepackage{booktabs}
\usepackage{array}
\usepackage{xcolor}
\usepackage{hyperref}
\usepackage{listings}
\usepackage{fancyhdr}
\usepackage{titlesec}
\usepackage{enumitem}
\usepackage{float}
\usepackage{caption}
\usepackage{mdframed}
\usepackage{lmodern}
\usepackage{microtype}
\usepackage{longtable}
\usepackage{multirow}
\usepackage{tabularx}
\usepackage{seqsplit}

\geometry{top=2.5cm,bottom=2.5cm,left=2.5cm,right=2.5cm,headheight=15pt}

\definecolor{anarenk}{RGB}{180,10,40}
\definecolor{morrenk}{RGB}{99,102,241}
\definecolor{yesil}{RGB}{74,222,128}
\definecolor{gri}{RGB}{120,130,140}
\definecolor{kodbg}{RGB}{26,31,46}
\definecolor{turuncu}{RGB}{245,158,11}
\definecolor{acikgri}{RGB}{245,245,245}

\hypersetup{
    colorlinks=true, linkcolor=anarenk,
    urlcolor=morrenk, citecolor=morrenk,
    pdftitle={CV-İlan Eşleşmesi: Final Ödev-2},
}

\lstdefinestyle{kod}{
    basicstyle=\ttfamily\small\color{white},
    keywordstyle=\color{morrenk}\bfseries,
    stringstyle=\color{yesil},
    commentstyle=\color{gri}\itshape,
    numberstyle=\tiny\color{gri},
    numbers=left, numbersep=8pt,
    backgroundcolor=\color{kodbg},
    frame=single, framerule=0.4pt,
    rulecolor=\color{gri},
    breaklines=true, tabsize=4,
    showstringspaces=false, xleftmargin=12pt,
}

\mdfdefinestyle{bilgi}{
    linecolor=anarenk, linewidth=1.5pt,
    backgroundcolor=acikgri,
    innertopmargin=8pt, innerbottommargin=8pt,
    innerleftmargin=10pt, innerrightmargin=10pt,
    roundcorner=4pt
}

\titleformat{\section}{\Large\bfseries\color{anarenk}}{\thesection}{1em}{}[\color{anarenk}\titlerule]
\titleformat{\subsection}{\large\bfseries\color{morrenk}}{\thesubsection}{1em}{}
\titleformat{\subsubsection}{\normalsize\bfseries}{\thesubsubsection}{1em}{}

\pagestyle{fancy}
\fancyhf{}
\fancyhead[L]{\small\color{gri}CV -- İlan Eşleşmesi | Proje [50]}
\fancyhead[R]{\small\color{gri}Doğal Dil İşleme $\cdot$ 2025--2026}
\fancyfoot[C]{\small\color{gri}\thepage}
\renewcommand{\headrulewidth}{0.4pt}

% ─────────────────────────────────────────────────────────────
\begin{document}
% ─────────────────────────────────────────────────────────────

% ── Kapak ────────────────────────────────────────────────────
\begin{titlepage}
\thispagestyle{empty}
\centering

{\large\bfseries T.C.}\\[0.2cm]
{\large\bfseries GÜMÜŞHANE ÜNİVERSİTESİ}\\[0.1cm]
{\large\bfseries GÜMÜŞHANE MESLEK YÜKSEKOKULU}\\[0.1cm]
{\large\bfseries BİLGİSAYAR TEKNOLOJİLERİ BÖLÜMÜ}\\[0.1cm]
{\large\bfseries BİLGİSAYAR PROGRAMCILIĞI PROGRAMI}\\[0.8cm]


\includegraphics[width=]{GU_LOGO.png}\\[0.3cm]

{\color{anarenk}\rule{0.85\textwidth}{1.5pt}}\\[0.4cm]

{\Large\bfseries\color{anarenk} FINAL ÖDEV-2}\\[0.2cm]
{\Large\bfseries\color{anarenk} CV -- İLAN EŞLEŞMESİ}\\[0.2cm]
{\large\bfseries\color{morrenk} Word2Vec Tabanlı Metin Benzerliği Hesaplama ve Değerlendirme}\\[0.2cm]
{\normalsize\itshape Proje [50]: İnsan Kaynakları -- Aday CV'leri ile İş İlanlarının\\
Metinsel Uyumunun Ölçülmesi}\\[0.3cm]

{\color{anarenk}\rule{0.85\textwidth}{1.5pt}}\\[0.8cm]

{\large\bfseries PROJE EKİBİ:}\\[0.4cm]
\begin{tabular}{rl}
    1. & \textbf{Mario Enrique Motede Dasilva | 2411081503} \\
    2. & \textbf{Heriberto Fernandez Chale | 2411081509} \\
    3. & \textbf{Matias Fernando Ndong Owono Obiang | 2311081634} \\
\end{tabular}

\vspace{0.8cm}
\begin{tabular}{ll}
    \textbf{Ders:} Doğal Dil İşleme \\[0.2cm]
    \textbf{Öğretim Elemanı:}   Dr. Rabia Yaşa Koştaş \\[0.2cm]
    \textbf{GitHub:} \url{https://github.com/herfech/CV-ilan-eslesmesi.git} \\
\end{tabular}

\vfill
{\color{anarenk}\bfseries 2025--2026 Bahar Dönemi}
\end{titlepage}

\tableofcontents
\newpage

% ─────────────────────────────────────────────────────────────
\section{Giriş}
% ─────────────────────────────────────────────────────────────

\subsection{Ödevin Amacı}

Bu ödevin amacı, CV--İlan Eşleşmesi projesine ait veri seti üzerinde Word2Vec modelleri eğiterek metinler arası benzerlik hesaplamaları yapmak ve modellerin karşılaştırmalı başarımını değerlendirmektir. Ödev-1'de gerçekleştirilen metin ön işleme adımlarının çıktıları (lemmatized ve stemmed veri setleri) bu ödevde girdi olarak kullanılmıştır.

\subsection{Veri Seti}

\begin{mdframed}[style=bilgi, rightmargin=-1cm]
\textbf{CV--İlan Eşleşmesi Veri Seti}\\[0.3cm]
\begin{tabular}{ll}
Kaynak       & Proje[50] kapsamında özgün olarak oluşturulmuştur(\texttt{create\_dataset.py})\\
Toplam belge & 280 \\
CV sayısı    & 250 (10 kategori, 25'er belge) \\
İş ilanı     & 30 (10 kategori, 3'er seviye: Junior/Mid/Senior) \\
Kategoriler  & Data Science, Software Engineering, Marketing, Finance, \\
             & HR, Healthcare, Education, Sales, Project Management, Cybersecurity \\
Format       & \texttt{document\_id} $|$ \texttt{content} \\
\end{tabular}
\end{mdframed}

\subsection{Ham Veri vs Ön İşlenmiş Veri Karşılaştırması}

\begin{table}[H]
\centering
\caption{Ham Veri ve Ön İşlenmiş Veri Yapısal Karşılaştırması}
\begin{tabular}{lccc}
\toprule
\textbf{Özellik} & \textbf{Ham Veri} & \textbf{Lemmatized} & \textbf{Stemmed} \\
\midrule
Belge sayısı    & 280    & 280    & 280    \\
Ort. kelime/belge & $\approx$105  & $\approx$68     & $\approx$65     \\
Stopword        & Var    & Yok    & Yok    \\
Büyük/küçük harf & Karma & Küçük  & Küçük  \\
Kelime formu    & Orijinal & Sözlük kökü & Kök/gövde \\
Dosya             & cv\_jobs\_raw.csv & lemmatized.csv & stemmed.csv \\
\bottomrule
\end{tabular}
\end{table}

\subsection{GitHub Reposu}

Tüm kodlar, modeller ve veri setleri GitHub reposunda mevcuttur:\\
\url{https://github.com/herfech/CV-ilan-eslesmesi}



% ─────────────────────────────────────────────────────────────
\section{Yöntem}
% ─────────────────────────────────────────────────────────────

\subsection{Word2Vec Vektörleştirme Detayları}

Gensim kutuphanesi kullanilarak asagidaki 8 parametre seti ile hem lemmatized hem de stemmed veri setleri icin toplamda \textbf{16 Word2Vec modeli} egitilmistir. Ortak parametreler: \texttt{min\_count=2}, \texttt{epochs=10}, \texttt{workers=4}, \texttt{seed=42}.

\begin{table}[H]
\centering
\caption{Word2Vec Model Parametreleri ve Tam Adlari}
\begin{tabular}{cccccl}
\toprule
\textbf{No} & \textbf{Veri Seti} & \textbf{Model Tipi} & \textbf{Window} & \textbf{Dim} & \textbf{Model Adi} \\
\midrule
1  & lemmatized & CBOW      & 2 & 100 & \texttt{word2vec\_lem\_cbow\_win2\_dim100} \\
2  & lemmatized & Skip-gram & 2 & 100 & \texttt{word2vec\_lem\_skipgram\_win2\_dim100} \\
3  & lemmatized & CBOW      & 4 & 100 & \texttt{word2vec\_lem\_cbow\_win4\_dim100} \\
4  & lemmatized & Skip-gram & 4 & 100 & \texttt{word2vec\_lem\_skipgram\_win4\_dim100} \\
5  & lemmatized & CBOW      & 2 & 300 & \texttt{word2vec\_lem\_cbow\_win2\_dim300} \\
6  & lemmatized & Skip-gram & 2 & 300 & \texttt{word2vec\_lem\_skipgram\_win2\_dim300} \\
7  & lemmatized & CBOW      & 4 & 300 & \texttt{word2vec\_lem\_cbow\_win4\_dim300} \\
8  & lemmatized & Skip-gram & 4 & 300 & \texttt{word2vec\_lem\_skipgram\_win4\_dim300} \\
9  & stemmed    & CBOW      & 2 & 100 & \texttt{word2vec\_stem\_cbow\_win2\_dim100} \\
10 & stemmed    & Skip-gram & 2 & 100 & \texttt{word2vec\_stem\_skipgram\_win2\_dim100} \\
11 & stemmed    & CBOW      & 4 & 100 & \texttt{word2vec\_stem\_cbow\_win4\_dim100} \\
12 & stemmed    & Skip-gram & 4 & 100 & \texttt{word2vec\_stem\_skipgram\_win4\_dim100} \\
13 & stemmed    & CBOW      & 2 & 300 & \texttt{word2vec\_stem\_cbow\_win2\_dim300} \\
14 & stemmed    & Skip-gram & 2 & 300 & \texttt{word2vec\_stem\_skipgram\_win2\_dim300} \\
15 & stemmed    & CBOW      & 4 & 300 & \texttt{word2vec\_stem\_cbow\_win4\_dim300} \\
16 & stemmed    & Skip-gram & 4 & 300 & \texttt{word2vec\_stem\_skipgram\_win4\_dim300} \\
\bottomrule
\end{tabular}
\end{table}

\subsection{Vektör Çıktıları — Her Model için En Benzer 5 Kelime}

Her model icin \texttt{``develop''} anahtar kelimesi secilerek vektör uzayindaki en yakin 5 kelime sunulmustur.
Amac, modelden modele benzerlik skorlarinin nasil degistigini gostermektir. Bu skorlar tek basina modelin
basarisini gostermek icin yetersizdir.

\begin{table}[H]
\centering
\caption{16 Model icin Vektor Ciktilari --- ``develop'' Kelimesi}
\small
\begin{tabular}{p{5.2cm}p{7cm}c}
\toprule
\textbf{Model Adi} & \textbf{En Benzer 5 Kelime} & \textbf{Maks. Skor} \\
\midrule
lem\_cbow\_win2\_dim100      & skill, design, implement, create, build         & 0.91 \\
lem\_skipgram\_win2\_dim100  & implement, build, create, deploy, maintain      & 0.89 \\
lem\_cbow\_win4\_dim100      & implement, design, build, create, deploy        & 0.92 \\
lem\_skipgram\_win4\_dim100  & implement, build, deploy, create, design        & 0.90 \\
lem\_cbow\_win2\_dim300      & implement, build, design, create, deploy        & 0.93 \\
lem\_skipgram\_win2\_dim300  & implement, build, deploy, maintain, create      & 0.91 \\
lem\_cbow\_win4\_dim300      & implement, build, design, create, deploy        & 0.94 \\
lem\_skipgram\_win4\_dim300  & implement, build, deploy, create, maintain      & 0.93 \\
stem\_cbow\_win2\_dim100     & skill, design, implement, creat, build          & 0.90 \\
stem\_skipgram\_win2\_dim100 & implement, build, creat, deploy, maintain       & 0.88 \\
stem\_cbow\_win4\_dim100     & implement, design, build, creat, deploy         & 0.91 \\
stem\_skipgram\_win4\_dim100 & implement, build, deploy, creat, design         & 0.89 \\
stem\_cbow\_win2\_dim300     & implement, build, design, creat, deploy         & 0.92 \\
stem\_skipgram\_win2\_dim300 & implement, build, deploy, maintain, creat       & 0.90 \\
stem\_cbow\_win4\_dim300     & implement, build, design, creat, deploy         & 0.93 \\
stem\_skipgram\_win4\_dim300 & implement, build, deploy, creat, maintain       & 0.92 \\
\bottomrule
\end{tabular}
\end{table}

\textit{Gozlem: dim=300 modellerinin maks. skoru dim=100 modellerine gore genel olarak daha yuksektir. Stemmed modellerde
``create'' yerine ``creat'', ``build'' yerine bazi durumlarda ``build'' kalir --- koke indirgeme anlamsal bilgiyi kismen kaybettirir.
Hangi modelin daha basarili olmasini bekliyoruz? Skip-gram + window=4 + dim=300 kombinasyonu: Skip-gram nadir kelimelerin
vektorlerini daha iyi ogrenirken genis pencere daha zengin baglamsal iliskiler saglar ve yuksek boyut daha ayirt edici
semantik uzaylar olusturur.}


\subsection{Ornek Giris Metni}

Asagidaki metin, veri setinden (\texttt{cv\_jobs\_dataset.csv}) secilmis olup disaridan alinmamistir.

\begin{mdframed}[style=bilgi]
\textbf{Document ID:} \texttt{cv\_001} \quad \textbf{Kategori:} Data Science\\[0.3cm]
\textit{Isabella Wilson principal professional with 18 years experience in data science field.
SKILLS AND EXPERTISE: python machine learning deep learning neural networks tensorflow pytorch
scikit-learn pandas numpy data analysis statistical modeling regression classification clustering
natural language processing computer vision data visualization matplotlib seaborn sql database
big data spark hadoop feature engineering model deployment mlops statistical analysis predictive
modeling data mining exploratory data analysis hypothesis testing ab testing time series forecasting
business intelligence tableau power bi r programming data warehousing etl pipeline cloud computing
aws gcp azure jupyter notebook git version control. PROFESSIONAL EXPERIENCE: developed machine
learning models to predict customer churn reducing attrition by 25 percent built recommendation
system increasing user engagement by 40 percent created automated data pipeline processing 10
million records daily. EDUCATION: master degree data science machine learning artificial intelligence
university. PROFESSIONAL SUMMARY: results driven Data Science professional with 18 years of
progressive experience.}
\end{mdframed}

\subsection{Benzerlik Hesaplama Yontemi}

Her model icin benzerlik asagidaki adimlarla hesaplanmistir:

\begin{enumerate}
    \item Giris metni ve karsilastirilacak belge tokenize edilir
    \item Her token icin Word2Vec modelinden vektor alinir; modelde olmayan kelimeler atlanir
    \item Tum token vektorlerinin aritmetik ortalamasi alinir (belge vektoru)
    \item Eger hic token vektoru yoksa Sifir Vektoru (Zero Vector) atanir --- NaN/ZeroDivisionError'u onler
    \item Giris metni ile her belge arasinda Kosinus Benzerligi hesaplanir
    \item Her model icin en yuksek skora sahip Top-5 belge secilir
\end{enumerate}

\begin{equation}
\cos(\theta) = \frac{\mathbf{A} \cdot \mathbf{B}}{\|\mathbf{A}\| \cdot \|\mathbf{B}\|}
\end{equation}

\begin{lstlisting}[style=kod, language=Python, caption={Belge Vektoru + Zero Vector Savunma Mekanizmasi}]
def doc_vector(text, model):
    tokens = preprocess(text)
    vecs = [model.wv[w] for w in tokens if w in model.wv]
    if not vecs:                          # Hic vektor yoksa
        return np.zeros(model.vector_size) # Zero Vector
    return np.mean(vecs, axis=0)
\end{lstlisting}


% ─────────────────────────────────────────────────────────────
\section{Sonuçlar ve Değerlendirme}
% ─────────────────────────────────────────────────────────────

\subsection{Değerlendirme-1: Cosine Benzerlik Tablosu}

Giris metnine (\texttt{cv\_001}) en benzer 5 belge her model icin cosine benzerligi ile bulunmustur.
Toplam: 16 model $\times$ 5 belge = 80 sonuc.

\begin{table}[H]
\centering
\caption{Her Model için Top-5 Benzer Belge ve Cosine Skorları (Tam Veri Seti)}
\label{tab:model_karsilastirma_tam}
\footnotesize % Hace la letra un punto más pequeña para que el texto largo quepa perfectamente
\begin{tabularx}{\textwidth}{l X X c}
\toprule
\textbf{Model Adı} & \textbf{5 Benzer Belge} & \textbf{Cosine Skorları} & \textbf{Ort.} \\
\midrule
lem\_cbow\_win2\_dim300 & cv\_136, cv\_234, cv\_004, cv\_074, cv\_154 & [1.0, 1.0, 1.0, 1.0, 1.0] & 0.9999 \\
lem\_cbow\_win4\_dim300 & cv\_234, cv\_136, cv\_226, cv\_096, cv\_154 & [1.0, 1.0, 1.0, 1.0, 1.0] & 0.9998 \\
stem\_cbow\_win2\_dim300 & cv\_011, cv\_121, cv\_044, cv\_074, cv\_004 & [1.0, 1.0, 1.0, 1.0, 1.0] & 0.9998 \\
lem\_cbow\_win2\_dim100 & cv\_136, cv\_234, cv\_096, cv\_146, cv\_226 & [1.0, 1.0, 1.0, 1.0, 1.0] & 0.9997 \\
stem\_cbow\_win2\_dim100 & cv\_011, cv\_121, cv\_073, cv\_100, cv\_211 & [1.0, 1.0, 1.0, 1.0, 1.0] & 0.9997 \\
stem\_cbow\_win4\_dim300 & cv\_011, cv\_073, cv\_121, cv\_154, cv\_234 & [1.0, 1.0, 1.0, 1.0, 1.0] & 0.9997 \\
lem\_cbow\_win4\_dim100 & cv\_136, cv\_096, cv\_226, cv\_196, cv\_176 & [1.0, 1.0, 1.0, 1.0, 1.0] & 0.9996 \\
stem\_cbow\_win4\_dim100 & cv\_011, cv\_121, cv\_073, cv\_211, cv\_234 & [1.0, 1.0, 1.0, 0.999, 0.999] & 0.9995 \\
stem\_skipgram\_win2\_dim300 & cv\_011, cv\_121, cv\_211, cv\_131, cv\_061 & [1.0, 0.999, 0.999, 0.997, 0.996] & 0.9983 \\
lem\_skipgram\_win2\_dim300 & cv\_121, cv\_011, cv\_211, cv\_172, cv\_140 & [0.999, 0.999, 0.999, 0.997, 0.997] & 0.9980 \\
stem\_skipgram\_win2\_dim100 & cv\_011, cv\_121, cv\_211, cv\_131, cv\_061 & [1.0, 1.0, 0.999, 0.996, 0.995] & 0.9978 \\
lem\_skipgram\_win2\_dim100 & cv\_121, cv\_011, cv\_211, cv\_172, cv\_031 & [0.999, 0.999, 0.999, 0.996, 0.995] & 0.9975 \\
stem\_skipgram\_win4\_dim300 & cv\_011, cv\_121, cv\_211, cv\_031, cv\_201 & [0.999, 0.999, 0.999, 0.992, 0.99] & 0.9957 \\
lem\_skipgram\_win4\_dim300 & cv\_011, cv\_121, cv\_211, cv\_031, cv\_061 & [0.999, 0.999, 0.999, 0.992, 0.989] & 0.9956 \\
stem\_skipgram\_win4\_dim100 & cv\_011, cv\_121, cv\_211, cv\_031, cv\_201 & [0.999, 0.999, 0.998, 0.991, 0.989] & 0.9952 \\
lem\_skipgram\_win4\_dim100 & cv\_011, cv\_121, cv\_211, cv\_031, cv\_201 & [0.999, 0.999, 0.999, 0.993, 0.987] & 0.9951 \\
\bottomrule
\end{tabularx}
\end{table}


\subsection{Değerlendirme-2: Anlamsal Değerlendirme}

Her modelin onerdigi 5 benzer belge icin 1--5 arasi anlamsal benzerlik puani \textbf{elle verilmistir} (YZ kullanilmamistir).
Giris metni Data Science CV'sidir; ayni kategorideki belgeler yuksek puan almistir.

\textbf{Puanlama Kriteri:}
\begin{itemize}
    \item \textbf{5 puan:} Neredeyse aynı temada, çok güçlü benzerlik
    \item \textbf{4 puan:} Anlamlı, açık benzerlik içeriyor
    \item \textbf{3 puan:} Ortalama düzeyde benzer
    \item \textbf{2 puan:} Kısmen ilgili ama bağlamı tutmuyor
    \item \textbf{1 puan:} Çok alakasız, anlamca zayıf benzerlik
\end{itemize}

\begin{table}[H]
\centering
\caption{16 Model --- Anlamsal Degerlendirme Puanlari (Elle Verilmistir)}
\footnotesize
\begin{tabularx}{\textwidth}{l X c c}
\toprule
\textbf{Model Adi} & \textbf{5 Benzer Belge} & \textbf{Puanlar (1--5)} & \textbf{Ort.} \\
\midrule
lem\_cbow\_win2\_dim300     & cv\_136, cv\_234, cv\_004, cv\_074, cv\_154 & 5, 5, 4, 4, 4 & 4.4 \\
lem\_cbow\_win4\_dim300     & cv\_234, cv\_136, cv\_226, cv\_096, cv\_154 & 5, 5, 4, 4, 3 & 4.2 \\
stem\_cbow\_win2\_dim300    & cv\_011, cv\_121, cv\_044, cv\_074, cv\_004 & 5, 5, 4, 4, 4 & 4.4 \\
lem\_cbow\_win2\_dim100     & cv\_136, cv\_234, cv\_096, cv\_146, cv\_226 & 5, 5, 4, 4, 4 & 4.4 \\
stem\_cbow\_win2\_dim100    & cv\_011, cv\_121, cv\_073, cv\_100, cv\_211 & 5, 5, 4, 3, 3 & 4.0 \\
stem\_cbow\_win4\_dim300    & cv\_011, cv\_073, cv\_121, cv\_154, cv\_234 & 5, 4, 5, 4, 3 & 4.2 \\
lem\_cbow\_win4\_dim100     & cv\_136, cv\_096, cv\_226, cv\_196, cv\_176 & 5, 5, 4, 4, 4 & 4.4 \\
stem\_cbow\_win4\_dim100    & cv\_011, cv\_121, cv\_073, cv\_211, cv\_234 & 5, 5, 4, 3, 3 & 4.0 \\
stem\_skipgram\_win2\_dim300 & cv\_011, cv\_121, cv\_211, cv\_131, cv\_061 & 5, 5, 4, 3, 3 & 4.0 \\
lem\_skipgram\_win2\_dim300 & cv\_121, cv\_011, cv\_211, cv\_172, cv\_140 & 5, 5, 4, 3, 3 & 4.0 \\
stem\_skipgram\_win2\_dim100 & cv\_011, cv\_121, cv\_211, cv\_131, cv\_061 & 5, 5, 4, 3, 3 & 4.0 \\
lem\_skipgram\_win2\_dim100 & cv\_121, cv\_011, cv\_211, cv\_172, cv\_031 & 5, 5, 4, 3, 2 & 3.8 \\
stem\_skipgram\_win4\_dim300 & cv\_011, cv\_121, cv\_211, cv\_031, cv\_201 & 5, 5, 4, 3, 3 & 4.0 \\
lem\_skipgram\_win4\_dim300 & cv\_011, cv\_121, cv\_211, cv\_031, cv\_061 & 5, 5, 4, 3, 3 & 4.0 \\
stem\_skipgram\_win4\_dim100 & cv\_011, cv\_121, cv\_211, cv\_031, cv\_201 & 5, 5, 4, 3, 3 & 4.0 \\
lem\_skipgram\_win4\_dim100 & cv\_011, cv\_121, cv\_211, cv\_031, cv\_201 & 5, 5, 4, 3, 2 & 3.8 \\
\bottomrule
\end{tabularx}
\end{table}

\textbf{Yorum:}
Neredeyse tum modeller yuksek anlamsal puan almistir (3.8--4.4). Bunun temel nedeni, giris metninin Data Science
kategorisinden olmasi ve modellerin cogunlugunun ayni kategoriden belgeler onermesidir. CBOW modellerinin anlamsal
puanlari Skip-gram modellerine kiyasla biraz daha yuksek cikabilmistir --- bu paradoks, CBOW'un sik gecen Data Science
kelimelerini daha iyi yakalamasindan kaynaklanmaktadir. Hangi modeller daha yuksek ortalama aldi? Genel olarak
\texttt{dim300} ve genis baglamli (\texttt{win4}) modeller; ancak bu veri seti icin farklar oldukca kucuktur.

\subsection{Değerlendirme-3: Jaccard Sıralama Tutarlılığı}

\[
\text{Jaccard}(A, B) = \frac{|A \cap B|}{|A \cup B|}
\]

Iki farkli modelin ayni giris metnine verdigi Top-5 sonuclar arasindaki ortusmeyi olcer. Tekil bir basari olcmutur; hangi modellerin birbirine benzer sonuclar urettigi gosterilir.

16$\times$16 Jaccard matrisi \texttt{data/jaccard\_matrix.csv} dosyasinda mevcuttur.

\begin{figure}[H]
\centering
\includegraphics[width=30\textwidth]{04_jaccard_heatmap.png}
\caption{Jaccard Benzerlik Matrisi Heatmap (16$\times$16) --- Modellerin Top-5 Sonuc Ortusme Orani.
Kosegen her zaman 1.00'dir (modelin kendisiyle karsilastirmasi); yorumlamada hesaba katilmaz.}
\end{figure}

\textbf{Temel bulgular:}
\begin{itemize}
    \item Ayni veri seti (lemmatized veya stemmed) ve benzer hiperparametreli modeller yuksek Jaccard skoru uretmistir
    \item CBOW modelleri kendi aralarinda, Skip-gram modelleri kendi aralarinda daha yuksek ortusme gostermistir --- mimari benzerlik siralama tutarliligini belirleyen temel faktordur
    \item window=4 modelleri, window=2 modellerine kiyasla birbirleriyle daha fazla orusen sonuclar uretmistir
    \item Lemmatized ve stemmed veri setleri arasindaki modeller gorece dusuk Jaccard skoru uretmistir; on isleme yontemi model ciktilarini onemli olcude etkilemektedir
    \item CBOW + dim300 grubu en tutarli i\c{c} grup olusturmaktadir
\end{itemize}

\subsection{Model Yapilandirmalarinin Etkisi}

\begin{table}[H]
\centering
\caption{Hiperparametre Etkisi Ozeti}
\begin{tabular}{lp{9cm}}
\toprule
\textbf{Parametre} & \textbf{Gozlem} \\
\midrule
CBOW vs Skip-gram & Skip-gram modelleri, ozellikle kucuk-orta veri setlerinde baglamsal iliskileri daha zengin ogrenirken bu veri setinde CBOW modelleri daha yuksek cosine skorlari uretmistir. \\
\midrule
Window boyutu (2 vs 4) & Window=4 modelleri daha genis baglamdan faydalanarak kavramsal iliskileri daha iyi ogrenmis, Jaccard matrisinde daha tutarli sonuclar uretmistir. \\
\midrule
Vektor boyutu (100 vs 300) & dim=300 modelleri semantik ozellikleri daha iyi yakalamistir; ancak egitim suresi dim=100'e kiyasla 2x uzamistir. \\
\midrule
Lemmatized vs Stemmed & Her iki on isleme yontemi de yuksek cosine skoru uretmistir. Stemming bazi kelimeleri kucultmekte (``create'' $\to$ ``creat'') ve anlam kaybina yol acabilmektedir. \\
\bottomrule
\end{tabular}
\end{table}

% ─────────────────────────────────────────────────────────────
\section{Sonuç ve Öneriler}
% ─────────────────────────────────────────────────────────────

\subsection{Genel Çıkarımlar}

Bu odevde CV--Ilan Eslesmesi projesine ait 280 belgeden olusan veri seti uzerinde 16 farkli Word2Vec modeli egitilmis ve 3 farkli degerlendirme yontemi uygulanmistir.

\textbf{En basarili modeller (cosine ortalama):}
\begin{enumerate}
    \item \texttt{word2vec\_lemmatized\_cbow\_win2\_dim300} --- Ortalama cosine: 0.9999
    \item \texttt{word2vec\_lemmatized\_cbow\_win4\_dim300} --- Ortalama cosine: 0.9998
    \item \texttt{word2vec\_stemmed\_cbow\_win2\_dim300} --- Ortalama cosine: 0.9998
\end{enumerate}

\textbf{En yuksek anlamsal puan:}
\begin{itemize}
    \item \texttt{\seqsplit{lem\_cbow\_win2\_dim300}}, \texttt{lem\_cbow\_win2\_dim100}, \texttt{lem\_cbow\_win4\_dim100}, \texttt{\seqsplit{stem\_cbow\_win2\_dim300}} --- Ortalama anlamsal puan: 4.4
\end{itemize}

\textbf{En dusuk basarimli model:}
\begin{itemize}
    \item \texttt{\seqsplit{word2vec\_lemmatized\_skipgram\_win4\_dim100}} ve \texttt{\seqsplit{word2vec\_stemmed\_skipgram\_win4\_dim100}} --- Anlamsal ortalama: 3.8
\end{itemize}

\subsection{Hangi Model Hangi Görev İçin?}

\begin{table}[H]
\centering
\caption{Görev Bazında Model Önerileri}
\begin{tabular}{ll}
\toprule
\textbf{Görev} & \textbf{Önerilen Model} \\
\midrule
Hızlı ön eleme (büyük veri) & CBOW + window2 + dim100 \\
Semantik doğruluk & Skip-gram + window4 + dim300 \\
Denge (hız + doğruluk) & Skip-gram + window2 + dim100 \\
Üretim ortamı & TF-IDF ön eleme + Word2Vec son sıralama \\
\bottomrule
\end{tabular}
\end{table}

\subsection{Gelecek Çalışmalar}

\begin{enumerate}
    \item \textbf{BERT / Sentence-Transformers:} Bağlam duyarlı modeller daha iyi anlam temsili sağlar
    \item \textbf{Daha büyük veri seti:} Gerçek CV ve iş ilanı verileriyle modelleri yeniden eğitmek
    \item \textbf{ATS Entegrasyonu:} Greenhouse, Workable, BambooHR API bağlantısı
    \item \textbf{PDF Desteği:} pdfplumber ile gerçek CV dosyalarını okuma
    \item \textbf{Streamlit Cloud:} Uygulamayı halka açık URL ile yayınlamak
\end{enumerate}

% ─────────────────────────────────────────────────────────────
\section*{Ekler}
\addcontentsline{toc}{section}{Ekler}
% ─────────────────────────────────────────────────────────────

\subsection*{A. Proje Dosya Yapısı}

\begin{lstlisting}[style=kod, language={}, caption={GitHub Repo Yapısı}]
CV-ilan-eslesmesi/
|-- app.py                    # Streamlit dashboard
|-- matching_engine.py        # TF-IDF + Cosine Similarity motoru
|-- sample_data.py            # 5 ornek CV ve is ilani
|-- create_dataset.py         # 280 belgeli veri seti olusturucu
|-- test_matching.py          # Terminal testi
|-- Final_Odev2.ipynb         # Odev-2 ana notebook (Word2Vec)
|-- data/
|   |-- cv_jobs_raw.csv       # Ham veri seti (280 belge)
|   |-- lemmatized.csv        # On islenmis (lemmatization)
|   |-- stemmed.csv           # On islenmis (stemming)
|   |-- cosine_evaluation.csv # Cosine degerlendirme sonuclari
|   |-- jaccard_matrix.csv    # 16x16 Jaccard matrisi
|-- models/                   # 16 Word2Vec modeli (.model)
|-- plots/                       # Grafik ciktilari (.png)
|-- notebooks/plots
|   |-- 01_eda_ve_onisleme.ipynb
|   |-- 02_model_karsilastirma.ipynb
|-- rapor/
    |-- cv_ilan_eslesme_rapor.tex
    |-- cv_ilan_eslesme_rapor.pdf
    |-- GU_LOGO.png
    |-- 04_jaccard_heatmap.png
\end{lstlisting}

\subsection*{B. Gereksinimler}

\begin{lstlisting}[style=kod, language={}, caption={requirements.txt}]
streamlit>=1.35.0
scikit-learn>=1.4.0
pandas>=2.0.0
plotly>=5.18.0
nltk>=3.8.0
numpy>=1.26.0
gensim>=4.3.0
pdfplumber>=0.10.0
matplotlib>=3.8.0
seaborn>=0.13.0
jupyter>=1.0.0
ipykernel>=6.0.0
\end{lstlisting}

\subsection*{C. Kaynaklar}

\begin{itemize}
    \item Mikolov, T. et al. (2013). Efficient Estimation of Word Representations in Vector Space. \textit{ICLR 2013}.
    \item Řehůřek, R., \& Sojka, P. (2010). Software Framework for Topic Modelling with Large Corpora. \textit{LREC 2010}.
    \item Gensim Documentation: \url{https://radimrehurek.com/gensim/}
    \item Scikit-learn Documentation: \url{https://scikit-learn.org}
    \item Streamlit Documentation: \url{https://docs.streamlit.io}
\end{itemize}

\end{document}
